\begin{abstract}
% area
Large language models (LLMs) are increasingly deployed as agents that interact with external software through tools, with an orchestration layer translating model outputs into structured tool invocations.
% problem
This architecture creates an LLM--orchestrator--tool trust boundary where tool semantics may go unenforced and LLM-generated arguments may be incorrectly treated as trusted.
% solution
We present \project{}, a tool-aware framework that derives semantic security checks for LLM tools from their existing metadata and a CWE-guided sanitizer corpus.
% methodology
\project{} specializes these checks to individual tools and automatically generates fuzzing harnesses and semantic oracles using runtime introspection and outcome sanitization.
% evaluation
We evaluated \project{} on 569 tools across 11 repositories and identified 14 previously unknown bugs corresponding to weakness classes including Arbitrary File Writing, MCP Server Crash, Incorrect Tool Description, and Denial of Service. Of these bugs, seven are input-independent and do not require a specific input to be triggered.
% takeaway
Our findings show that real-world tools frequently fail to enforce their documented input semantics, making it unsafe to treat LLM-generated arguments as trusted.

\end{abstract}